\documentclass[11pt]{article}
\usepackage[T1]{fontenc}
\usepackage{lmodern}
\usepackage[margin=1in]{geometry}
\usepackage{amsmath,amssymb,amsthm}
\usepackage{booktabs}
\usepackage[colorlinks=true,linkcolor=blue,citecolor=blue,urlcolor=blue]{hyperref}
\hypersetup{pdftitle={Maximizing the adjacent comparison classes of parking functions},pdfauthor={Qiyuan Ma},pdfsubject={Research preprint: parking function comparison encoding maximum}}
\setlength{\emergencystretch}{3em}
\newtheorem{theorem}{Theorem}[section]
\newtheorem{lemma}[theorem]{Lemma}
\newtheorem{proposition}[theorem]{Proposition}
\newtheorem{corollary}[theorem]{Corollary}
\theoremstyle{definition}
\newtheorem{definition}[theorem]{Definition}
\newcommand{\PF}{\operatorname{PF}}
\newcommand{\enc}{\operatorname{enc}}
\newcommand{\rev}{\operatorname{rev}}
\newcommand{\irev}{\operatorname{irev}}
\newcommand{\A}{\mathrm A}
\newcommand{\D}{\mathrm D}
\newcommand{\T}{\mathrm T}
\newcommand{\eps}{\varepsilon}
\title{Maximizing the adjacent comparison classes of parking functions}
\author{Qiyuan Ma\\
  \small Shanghai Jiao Tong University\\
  \small \href{mailto:mqyqwerty@sjtu.edu.cn}{mqyqwerty@sjtu.edu.cn}\\
  \small ORCID: \href{https://orcid.org/0009-0000-9532-0821}{0009-0000-9532-0821}}
\date{October 28, 2026\\ \small Research preprint, version 1.0.0}

\begin{document}
\maketitle
\begin{abstract}
A parking function of length $n$ records each adjacent comparison by
$\A$, $\D$, or $\T$, according as it is a strict ascent, a strict descent,
or a tie. We determine all comparison words with a largest class.
For odd $n\ge3$, the unique largest class has encoding
$\A\D\A\D\cdots\A\D$. For even $n\ge4$, the two alternating
encodings tie for the maximum. At $n=2$, all three one-letter classes
have size one. The proof uses explicit injections that strictly eliminate
tied candidates and repair the first pair of consecutive comparisons
in the same direction. We also give an exact subset recurrence for the
maximum and a reproducible collection of checks. This answers
Problem~54 of Cruz et al.\ in the form of a self-contained research
preprint; no independent expert peer review or historical-priority
certification is claimed.
\end{abstract}

\section{The maximizing words}
A word $a=(a_1,\ldots,a_n)\in[n]^n$ is a parking function if its
nondecreasing rearrangement $b_1\le\cdots\le b_n$ satisfies $b_i\le i$.
Its comparison encoding $\enc(a)=w_1\cdots w_{n-1}$ is defined by
\[
w_i=\A\quad(a_i<a_{i+1}),\qquad
w_i=\D\quad(a_i>a_{i+1}),\qquad
w_i=\T\quad(a_i=a_{i+1}).
\]
Write $F_n(w)=|\{a\in\PF_n:\enc(a)=w\}|$.
Let $Z_n^\A=\A\D\A\D\cdots$ and $Z_n^\D=\D\A\D\A\cdots$,
both of length $n-1$, and set $Z_1^\A=\eps$.

\begin{theorem}\label{thm:main}
All maximizing encodings are given by
\[
\operatorname*{arg\,max}_{w\in\{\A,\D,\T\}^{n-1}}F_n(w)=
\begin{cases}
\{\eps\},&n=1,\\
\{\A,\D,\T\},&n=2,\\
\{Z_n^\A\},&n\ge3\text{ odd},\\
\{Z_n^\A,Z_n^\D\},&n\ge4\text{ even}.
\end{cases}
\]
Consequently the maximum is $M_n=F_n(Z_n^\A)$ for every $n\ge1$.
The recurrence in Section~\ref{sec:recurrence} computes each $M_n$
without any remaining maximization.
\end{theorem}

Cruz et al.~\cite{Cruz2024} ask which full statistic encoding has
the largest class. The analogous extremal result for ordinary
permutation descent classes is classical; see, for example,
Sagan--Yeh--Ziegler~\cite{Sagan1994}. Repeated values prevent
an immediate transfer of a multiset-preserving permutation argument.
The alternating parking-function counts themselves were recorded in
2015 as OEIS entries A260694 and A264963~\cite{OEISup,OEISdown}.
We do not claim these numerical sequences as new.
Recent recurrences for aggregate peak and tie statistics
\cite{Alexandersson2026} concern a different refinement from a fixed
complete comparison word.

\section{Closure and rank assignments}\label{sec:closure}
For a positive-integer word $a$ define
$C_a(t)=|\{i:a_i\le t\}|$.
The parking condition is equivalent to
\[
C_a(t)\ge t\qquad(1\le t\le n).
\]

\begin{lemma}\label{lem:closure}
Reordering a parking function preserves the parking condition.
Replacing one entry by a smaller positive integer also preserves it.
If a parking function contains at least two copies of $x$, increasing
one copy to $x+1$ preserves the parking condition.
\end{lemma}
\begin{proof}
Reordering leaves every $C_a(t)$ unchanged, and decreasing an entry
can only increase these counts. For the last statement,
\[
C_a(x)\ge C_a(x-1)+2\ge x+1,
\]
where $C_a(0)=0$. The modification decreases only $C_a(x)$, by one;
all other integer-threshold counts are unchanged. Thus the inequalities
still hold. In particular, the modification cannot produce a value
larger than $n$ in a parking function.
\end{proof}

Collapse each maximal run of ties of a comparison word into one block.
The blocks form a directed path, with an arrow from the smaller block
to the larger block. Its transitive closure is a poset. Assigning
the distinct values $1,\ldots,m$ to its $m$ blocks in a linear extension
gives a parking function: each of the preceding $j-1$ block values
occurs at least once before the first occurrence of value $j$ in sorted
order. Thus every occurrence of $j$ has sorted position at least $j$.
We use this rank-assignment observation also after merging
nonadjacent minimal vertices or nonadjacent maximal vertices.
Such a merge does not create a directed cycle.

\section{The two odd alternating directions}\label{sec:odd}
The following map will also be used on prefixes.

\begin{lemma}\label{lem:R}
Let $k\ge3$ be odd, and suppose
$a_1>a_2<a_3>\cdots<a_k$. Define
\[
R(a)=
\begin{cases}
(a_2,\ldots,a_k,a_1),&a_1<a_k,\\
(a_k,a_1,\ldots,a_{k-1}),&a_1>a_k,\\
(a_2,\ldots,a_{k-1},a_1,a_2),&a_1=a_k.
\end{cases}
\]
Then $R$ is an injection into the alternating words with low
endpoints. Its last entry is strictly smaller than $a_k$.
Its multiset is unchanged, or one value is decreased.
\end{lemma}
\begin{proof}
The adjacent inequalities in each displayed tuple give the required
alternation. In the first two cases the last entries are $a_1<a_k$
and $a_{k-1}<a_k$. In the last case it is $a_2<a_1=a_k$;
one copy of $a_1$ is replaced by the smaller value $a_2$.
If the output is $b$, its endpoints satisfy respectively
$b_1<b_k$, $b_1>b_k$, and $b_1=b_k$.
Thus the branch is recognizable. The first two inverse maps are
the opposite cyclic shifts; the third inverse is
\[
(b_{k-1},b_1,\ldots,b_{k-2},b_{k-1}).
\]
\end{proof}

\begin{proposition}\label{prop:directions}
For odd $n\ge3$, $F_n(Z_n^\D)<F_n(Z_n^\A)$.
For even $n$, the two alternating classes have equal size.
\end{proposition}
\begin{proof}
Lemma~\ref{lem:R} and Lemma~\ref{lem:closure} give the injection
for odd $n$. It is not onto. The word
\[
b=(1,2)^{(n-3)/2}(1,n,1)
\]
is a parking function with encoding $Z_n^\A$; juxtaposition here
denotes concatenation. Its endpoints agree. Its only possible inverse
under $R$ contains two copies of $n$, so is not a parking function.
For even $n$, reversing all coordinates interchanges the two
alternating encodings and preserves the parking condition.
\end{proof}

\section{Strict repair of the first defect}\label{sec:binary}
For a binary comparison word $w\in\{\A,\D\}^{n-1}$, call
$w_i=w_{i+1}$ a defect. Suppose $w$ has a defect, and let
$k\in\{2,\ldots,n-1\}$ be the smallest index with
$w_{k-1}=w_k=c$. The first $k-1$ comparisons are alternating.
Let $w^+$ interchange $\A$ and $\D$ in these first $k-1$
positions and leave the other positions unchanged.

\begin{theorem}\label{thm:first-defect}
For this first-defect operation, $F_n(w)<F_n(w^+)$.
\end{theorem}

Set $s=(a_1,\ldots,a_k)$ and $t=a_{k+1}\ge1$.
All maps below leave $a_{k+1},\ldots,a_n$ fixed. Put
\[
\rev(s)=(a_k,\ldots,a_1),\qquad
\irev(s)=(a_1,a_{k-1},a_{k-2},\ldots,a_2,a_k),
\]
where the second map reverses only the interior and is the identity
when $k=2$.

\subsection{Even prefixes}
If $k$ is even and $c=\A$, the prefix starts low and ends high,
with $a_k<t$. Define
\[
E_\A(s;t)=
\begin{cases}
\rev(s),&a_1<t,\\
\irev(s),&a_1\ge t.
\end{cases}
\]
The first branch reverses the alternating direction and has new last
entry $a_1<t$. For the second branch,
\[
a_1\ge t>a_k>a_{k-1},\qquad
a_2>a_1\ge t>a_k.
\]
These inequalities check the new endpoint comparisons; the interior
comparisons reverse by alternation. The last entry is still $a_k<t$.
Thus both branches give the required encoding and preserve the
comparison with $t$. Their first output entries are respectively
less than $t$ and at least $t$. This distinguishes the branches;
the same reversal in each branch recovers the input.

If $k$ is even and $c=\D$, use
\[
E_\D(s;t)=
\begin{cases}
\rev(s),&a_1>t,\\
\irev(s),&a_1\le t.
\end{cases}
\]
Now $a_k>t$. In the second branch
\[
a_1\le t<a_k<a_{k-1},\qquad
a_2<a_1\le t<a_k,
\]
so the new endpoint comparisons are correct and the last entry
remains greater than $t$. The first output entries are respectively
greater than $t$ and at most $t$, giving the inverse branches.
All these operations merely reorder the prefix, hence preserve
the parking condition.

\subsection{Odd prefixes}
If $k$ is odd and $c=\A$, the prefix has high endpoints and
$a_k<t$. Apply $R$ from Lemma~\ref{lem:R}. It reverses the
alternating direction and decreases the last entry, preserving
the following ascent and the parking condition.

For the remaining case we need an upper-threshold version.

\begin{lemma}\label{lem:S}
Let $k\ge3$ be odd, let $a$ be alternating with low endpoints,
and suppose $a_k>t\ge1$. Then
\[
S_t(a)=
\begin{cases}
(a_k,a_1,\ldots,a_{k-1}),&a_1<a_k,\\
(a_2,\ldots,a_k,a_1),&a_1>a_k,\\
(a_2,\ldots,a_{k-1},t,a_k),&a_1=a_k
\end{cases}
\]
is an injection into high-endpoint alternating words, with last
entry at least $a_k>t$. It preserves the multiset or decreases
one value.
\end{lemma}
\begin{proof}
The first two branches are cyclic shifts; their last entries are
strictly greater than $a_k$. In the last branch write
$x=a_1=a_k>t$. The inequalities $a_2>x$ and $a_{k-1}>x>t$
give the required high-endpoint alternation; one copy of $x$
is replaced by the positive value $t$, and the last entry remains $x$.

Every output $b$ has unequal endpoints, so $R(b)$ uses only a
cyclic shift. On the first two branches $R(b)=a$, whose last entry
is greater than $t$. On the last branch
\[
R(b)=(x,a_2,\ldots,a_{k-1},t)
\]
ends in exactly $t$. This distinguishes the last branch from the
others. Its inverse is
$(b_k,b_1,\ldots,b_{k-2},b_k)$; the other inverse is $R(b)$.
\end{proof}

If $k$ is odd and $c=\D$, the prefix has low endpoints and
$a_k>t$. Apply $S_t$ to it. This preserves the following descent,
the required target encoding, and the parking condition by
Lemma~\ref{lem:closure}. We have constructed injections in all four
cases.

\subsection{Strictness in all four cases}
We exhibit a target parking function outside the image in each case.
All the target encodings here are binary.

For even $k$ and $c=\A$, the first vertex of the target prefix is
maximal in the full comparison poset, and its $k$th vertex is minimal.
Give them values $n$ and 1, and give the other vertices the distinct
values $2,\ldots,n-1$ in a linear extension. This permutation is a
parking function. Its first entry is at least $t$, so its only possible
inverse reverses the interior and retains first entry $n$. The
inverse cannot start with an ascent. It is not a source word.
For even $k$ and $c=\D$, give the target's first vertex value 1 and
its $k$th vertex value $n$, and rank the remaining vertices.
Its only possible inverse retains first entry 1 and cannot start
with a descent. This also covers $k=2$.

For odd $k$ and $c=\A$, the target prefix has low endpoints,
both minimal in the full comparison poset. Merge these endpoints
as one block of value 1, give the last prefix peak (vertex $k-1$)
value $n$, and give the other vertices the distinct ranks
$2,\ldots,n-2$. The sorted word is
$(1,1,2,\ldots,n-2,n)$, hence is a parking function.
Its only possible inverse under $R$ contains two copies of $n$.

For odd $k$ and $c=\D$, the target prefix has high endpoints,
both maximal in the full comparison poset. Merge them and rank
the resulting $n-1$ blocks. This gives a target parking function
whose prefix endpoints agree. All outputs of $S_t$ have unequal
prefix endpoints. This proves Theorem~\ref{thm:first-defect}.

\begin{corollary}\label{cor:binary}
Every nonalternating binary encoding is strictly smaller than some
alternating encoding.
\end{corollary}
\begin{proof}
Write $\delta(w)=|\{i:w_i=w_{i+1}\}|$. Flipping the first $k-1$
letters preserves all comparisons between letters within that
prefix, fixes the suffix, and changes the boundary pair
$w_{k-1}=w_k$ to an unequal pair. Thus
$\delta(w^+)=\delta(w)-1$. Iterating
Theorem~\ref{thm:first-defect} terminates at an alternating word.
\end{proof}

\section{Strict elimination of ties}\label{sec:ties}
\begin{theorem}\label{thm:ties}
For $n\ge3$, every encoding containing $\T$ is strictly smaller
than some binary encoding of the same length.
\end{theorem}
\begin{proof}
First construct non-strict injections which remove one tie at a time.
Choose the rightmost $\T$. If it is the last comparison, change
the right member of the pair $(x,x)$ to $x+1$, changing $\T$ to $\A$.
If its next letter is $\D$, the local triple is $(x,x,z)$ with
$x>z$; use $(x,x+1,z)$, changing $\T\D$ to $\A\D$.
Lemma~\ref{lem:closure} applies and subtracting the increment is
an inverse.

Otherwise the next comparison is $\A$. Let $q$ be the longest
alternating segment starting there, and let
$s=(s_1,\ldots,s_L)$ have encoding $\T q$, with $s_1=s_2=x$.
Any comparison after this segment is the same as the last letter of $q$,
because no tie lies to its right. The following maps keep the first
entry $x$ fixed.

If $q$ ends in $\A$, so $L$ is odd, use
\[
\Phi(s)=(x,s_3,\ldots,s_L,\min(x,s_{L-1})).
\]
The encoding is $q\D$. The new last entry is smaller than $s_L$,
so any following ascent remains valid. One copy of $x$ is replaced
by a value at most $x$. If $b$ is the output, the inverse is
$(b_1,b_1,b_2,\ldots,b_{L-1})$.

If $q$ ends in $\D$, so $L$ is even, put $u=s_L$ and use
\[
\Psi(s)=
\begin{cases}
(x,s_3,\ldots,s_L,x),&u<x,\\
(x,u,x,s_3,\ldots,s_{L-1}),&u>x,\\
(x,s_3,\ldots,s_{L-1},x,x+1),&u=x.
\end{cases}
\]
The encoding is $q\A$. The new last entries are respectively
$x>u$, $s_{L-1}>u$, and $x+1>u$, preserving any following descent.
The first two cases reorder; in the third, the input contains three
copies of $x$ and one is increased to $x+1$, as allowed by
Lemma~\ref{lem:closure}.
The output branches are distinguished by
\[
\begin{array}{c|c}
u<x& b_L=b_1=x\\
u>x& b_L>x,\quad b_2<b_L\\
u=x& b_L>x,\quad b_2\ge b_L .
\end{array}
\]
The corresponding inverse tuples are
\[
\begin{aligned}
&(x,x,b_2,\ldots,b_{L-1}),\\
&(x,x,b_4,\ldots,b_L,b_2),\\
&(x,x,b_2,\ldots,b_{L-2},x).
\end{aligned}
\]
All outside comparisons have been preserved. The target encoding
depends only on the source encoding, and exactly one $\T$ disappears.
Iteration therefore gives an injection to a fixed binary class.

We now check strictness. When the selected tie is followed by $\A$,
every image under $\Phi$ or $\Psi$ has a repeated value within the
new strictly alternating segment. Rank the target's tie blocks
distinctly as in Section~\ref{sec:closure}. The resulting parking
function has distinct values within that segment and is missing
from the image, so this step is strict.

If the selected tie is followed by $\D$, the image of the simple
increment has difference exactly 1 on its newly created ascent.
The target has a singleton peak at that ascent. Give this peak
value $n$. If the target has $m<n$ tie blocks, give the remaining
blocks the ranks $1,\ldots,m-1$, so its left neighbor is at most $n-2$.
If $m=n$, deleting the peak separates a nonempty left part from
a nonempty right part. Take a maximal element of the right part
last among the other vertices, with rank $n-1$. Again the peak's
left neighbor is at most $n-2$. These assignments are parking
functions and have a gap at least 2, so this step is strict.

It remains to consider a rightmost tie at the end. Apply the terminal
increment first. If another tie remains, its next letter is $\A$
or $\D$, and the next step is strict by the preceding arguments.
If there was just one tie, its preceding letter exists when $n\ge3$.
If that letter is $\D$, the target ends in $\D\A$.
Assign its last valley value 1, its final entry value $n$, and rank
the other vertices in between. The last gap is $n-1>1$, so this
target is not an image of the increment.
If the preceding letter is $\A$, the first target ends in $\A\A$.
Swap its last two entries, giving an injection to a word ending in
$\A\D$. Its image additionally requires the first of those last
three entries to be smaller than the last entry. Give the new
peak value $n$, the last entry value 1, and rank the other vertices
in between to obtain a missing target. Thus a strict comparison
is obtained in every case, proving the theorem.
\end{proof}

\section{Conclusion of the classification}
For $n\ge3$, Theorem~\ref{thm:ties} excludes every tied word.
Corollary~\ref{cor:binary} excludes every nonalternating binary word.
Proposition~\ref{prop:directions} selects $Z_n^\A$ for odd $n$
and gives equality between the two alternating words for even $n$.
All exclusions are strict. At $n=1$ there is only $(1)$, and at
$n=2$ the parking functions are $(1,1),(1,2),(2,1)$.
This proves Theorem~\ref{thm:main}, including all equality cases.

\section{An exact recurrence for the maximum}\label{sec:recurrence}
Let $Q_n$ be the poset on $V=[n]$ with
\[
1<_{Q_n}2>_{Q_n}3<_{Q_n}4>_{Q_n}\cdots .
\]
For a set $S$ of already assigned positions, let
\[
\mathcal M(S)=
\{i\notin S:i\text{ odd}\}
\ \cup\
\{i\notin S:i\text{ even and all adjacent odd positions lie in }S\}.
\]
Initialize $d_0(\varnothing)=1$ and $d_0(S)=0$ otherwise.
For $j=1,\ldots,n$, start a zero layer and perform, for every $S$
and every $R\subseteq\mathcal M(S)$,
\[
d_j(S\cup R)\mathrel{+}=d_{j-1}(S)
\qquad\text{if } |S\cup R|\ge j.
\]
Empty $R$ is allowed. Then
\[
\boxed{M_n=d_n([n]).}
\]
Indeed, $S$ records the positions with value at most $j-1$,
and $R$ records exactly those with value $j$. Strict order
constraints permit precisely subsets of the remaining minimal
elements. The cumulative-size inequality is exactly the parking
condition. Every permissible assignment is counted once.
An on-the-fly implementation has a crude time bound $O(n3^n)$
and two-layer space bound $O(2^n)$; the supplied implementation
caches transitions and uses arbitrary-precision integers.

\begin{center}
\begin{tabular}{r r@{\qquad}r r}
\toprule
$n$&$M_n$&$n$&$M_n$\\
\midrule
1&1&7&9470\\
2&1&8&104436\\
3&4&9&1432713\\
4&17&10&20709209\\
5&120&11&354493902\\
6&915&12&6343919118\\
\bottomrule
\end{tabular}
\end{center}

\section{Executable checks and provenance}
The repository provides the displayed injections, inverse maps,
missing-image witnesses, full ternary encoding counts through $n=10$,
and maximum values through $n=20$.
All $18\,248$ parking functions of lengths at most six were passed
through the complete composed injection and recovered by its inverse,
with injectivity checked separately for every source encoding.
Direct enumeration, weighted ordered set partitions, and poset
assignment agree on every class through $n=6$.
Every nonalternating binary encoding through $n=10$ was checked
against the strict first-defect inequality and a constructed
missing-image witness. These finite tests audit the implementation;
the proofs above do not depend on an enumeration cutoff.

\paragraph{AI assistance.}
OpenAI Codex assisted extensively with mathematical exploration,
proof development, programming, checking, literature screening,
and English exposition. This is an AI-assisted research preprint.
Independent expert peer review and proof-assistant formalisation
are not claimed. See the repository's \texttt{AI\_ASSISTANCE.md}
and \texttt{VERIFICATION.md} for the released scope.

\paragraph{Availability and licensing.}
The manuscript, documentation, and JSON data are licensed under
CC BY 4.0. The Python and C++ code is licensed under MIT.
The repository is
\url{https://github.com/KevinMaSAI/parking-function-encoding-maxima}.

\begin{thebibliography}{9}
\bibitem{Cruz2024}
A.~Cruz, P.~E.~Harris, K.~J.~Harry, J.~Kretschmann,
M.~McClinton, A.~Moon, J.~O.~Museus, and E.~Redmon,
\emph{On Some Discrete Statistics of Parking Functions},
Journal of Integer Sequences \textbf{27} (2024), Article 24.8.6.
\url{https://cs.uwaterloo.ca/journals/JIS/VOL27/Harris/harris13.pdf}.

\bibitem{Sagan1994}
B.~E.~Sagan, Y.-N.~Yeh, and G.~M.~Ziegler,
\emph{Maximizing M\"obius functions on subsets of Boolean algebras},
Discrete Mathematics \textbf{126} (1994), 293--311.
\url{https://doi.org/10.1016/0012-365X(94)90273-9}.

\bibitem{OEISup}
The OEIS Foundation Inc.,
\emph{A260694: Number of up-down parking functions of length $n$},
entry initiated in 2015.
\url{https://oeis.org/A260694}.

\bibitem{OEISdown}
The OEIS Foundation Inc.,
\emph{A264963: Number of down-up parking functions of length $n$},
entry initiated in 2015.
\url{https://oeis.org/A264963}.

\bibitem{Alexandersson2026}
P.~Alexandersson,
\emph{Parking functions, Smirnov words, and noncrossing Chow polynomials},
arXiv:2609.05131v1 (2026).
\url{https://arxiv.org/abs/2609.05131}.
\end{thebibliography}
\end{document}
